\begin{proposition}[KT-JACOBI-THIRD-UNIT-NONIMPLICATION]
A field arising from a unit heat coefficient in
KT-JACOBI-TWO-UNIT-FIELD need not have the third required unit,
even in the automatic family (JF.13).
\end{proposition}
\begin{proof}
Take $p=31$, $e=16=(p+1)/2$ and $D=-79$.  This is an odd
fundamental discriminant prime to $p$.  The exact character sums give
\[
\begin{gathered}
 B_{1,\chi_D}=-5,\qquad
 \frac{B_{3,\chi_D}}3=992=31\cdot32,\\
 \frac{B_{17,\chi_D}}{17}
       =-23090968227437542810000061098240
       \equiv20\pmod{31}.
\end{gathered}
\]
The reproducible standard-library certificate
\texttt{jacobi\_third\_unit\_counterexample.py} evaluates all three
numbers by the Bernoulli-polynomial formula and independently
reconciles the last residue using the defining conductor-$79\cdot31$
sum modulo $31^2$.  In particular $h(-79)=5$.
At $(n,r)=(20,1)$ the heat coefficient in (JF.11) is
$79(-20)\equiv1\pmod{31}$.  Also $K=50$, so $n\leq K/2$
and $|D|=79\leq3p+6=99$: this is inside the proved selection bound.
But $2e\equiv2\pmod{30}$, so (FV.2) identifies the third unit
with $B_{3,\chi_D}/3$, which is zero modulo $31$.
This disproves the automatic implication for an individually selected
field, not the existence of another field with all three units.
It is a scope discriminator and not a new rational vanishing result.
\end{proof}
